\chapter{MongoDB Setup}

\section{Core Concepts}

\begin{itemize}
  \item \textbf{MongoDB} --- document-oriented NoSQL database; stores flexible JSON-like documents.
  \item \textbf{MongoDB Atlas} --- MongoDB's managed cloud hosting.
  \item \textbf{Database} --- logical container for collections (e.g. \texttt{taskmanager}).
  \item \textbf{Collection} --- group of documents, analogous to a SQL table.
  \item \textbf{Document} --- a single BSON record, analogous to a row.
  \item \textbf{ObjectId} --- MongoDB's 12-byte unique id, auto-assigned to \texttt{\_id}.
  \item \textbf{BSON} --- binary-encoded storage format (adds types like Date and ObjectId).
\end{itemize}

\begin{diagrambox}[MongoDB Hierarchy]
\begin{verbatim}
Server/Cluster -> Database ("taskmanager") -> Collection ("tasks")
              -> Document ({ _id, title, status }) -> Fields
\end{verbatim}
\end{diagrambox}

\section{Connecting with Mongoose}

\whatwhyhow{
Mongoose is an ODM that sits between Node.js and MongoDB, giving schemas, validation, and query helpers.
}{
Raw driver calls are untyped and repetitive. Mongoose lets you define a schema once and get consistent validation everywhere.
}{
Create a single connection module (\texttt{config/db.js}) that reads the connection string from env vars and connects once at startup.
}

\subsection{config/db.js}

\begin{lstlisting}[language=JavaScript]
const mongoose = require("mongoose");

const connectDB = async () => {
  try {
    const conn = await mongoose.connect(process.env.MONGO_URI);
    console.log(`MongoDB connected: ${conn.connection.host}`);
  } catch (error) {
    console.error(`MongoDB connection error: ${error.message}`);
    process.exit(1);
  }
};

module.exports = connectDB;
\end{lstlisting}

\section{MongoDB Atlas Connection String}

\begin{lstlisting}[language=JavaScript]
// .env
MONGO_URI=mongodb+srv://<username>:<password>@cluster0.mongodb.net/taskmanager?retryWrites=true&w=majority
\end{lstlisting}

\begin{notebox}[NOTE]
Atlas requires whitelisting IPs under Network Access; use your current IP, or \texttt{0.0.0.0/0} temporarily for local testing only.
\end{notebox}

\begin{warnbox}[WARNING]
Never commit \texttt{.env} or hardcode credentials. Add \texttt{.env} to \texttt{.gitignore} and provide a \texttt{.env.example} placeholder.
\end{warnbox}
